\section*{译者说明}

这份 PDF 是 Sam Raskin 教授 2016 年春季在 MIT 讲授的《数论 II（18.786）》讲义的中文试译。底本为 Oron Propp 整理的英文讲义（2017 年 8 月 21 日版，共 23 讲，119 页）。译文由自动翻译流程生成、经人工合稿校对，用于检验「数学长文 + 检索插件」这条流水线；不保证与英文原稿逐字对应，引用请以英文底本为准。

几点约定：
\begin{itemize}[leftmargin=2em]
  \item 原讲义中的定义、定理、例、式均保留原有编号（如「定理 1.5」「式 (1.1)」），不重新编号；讲标题按原文顺序排列。
  \item 原稿是从 PDF 文本层提取的，公式在提取时普遍出现上下标错位、符号粘连。译文中的公式均按数学含义重排为 \LaTeX{}，并已逐条对照公开的原讲义 PDF（MIT OCW 版）核校；原文本身的排印错误，凡能确证者已在译文中订正，并在 \texttt{\% [原文疑误：…]} 注释中留痕（注释不进入 PDF），凡仅属抽取丢字、原文可通者则照原文排并加注说明。
  \item 术语以随书术语表为准（见文末）；人名首次出现时附英文原文，术语中的姓氏保留拉丁拼写（如 Artin 互反律、Tate 上同调）。
  \item 参考文献条目保留英文原貌，以便读者按原书索引查证。
\end{itemize}

\section*{前言}

这份讲义对应 2016 年春季学期由 Sam Raskin 在 MIT 讲授的《数论 II》（18.786）课程。课程原网页见 \url{http://math.mit.edu/~sraskin/cft/index.html}；除本讲义外，网页上还有带注释的参考书目与习题集，讲义中会经常引用。习题集是课程的重要部分，其中常常引入讲义正文没有覆盖的材料（不过索引里大多列了出来）。

我写这份讲义时，总体上尽量复现 Sam 讲课的内容，偶尔补全细节、或为便于自己理解而稍作重组。手写笔记在学期中就开始录入，约一半工作是在大二暑假有更多空闲时完成的。因此，讲义有些地方比较粗糙，但我猜越往后质量越好。当然，其中的错误都只属于我，与 Sam 无关；如有笔误、错误或其它意见，欢迎发邮件至 opropp@mit.edu。

课程从第 2–5 讲开始，用 Hilbert 符号分析类域论的二次情形，以便对类域论的思想作更具体的引入。第 6–13 讲发展此后处理类域论所需的上同调工具。接着，第 14–18 讲证明局部类域论的主定理；第 16 讲由 Marc Hoyois 在 Sam 缺席时代讲，谨此致谢。最后，第 19–23 讲证明整体类域论的主定理，这大概是全课最难的部分。课程以 Bjorn Poonen 教授的两讲收尾：第一讲关于二次型与 Hasse--Minkowski 定理，第二讲关于 Witt 消去与 Brauer--Manin 障碍；遗憾的是，我还没来得及把它们整理成文。

感谢 Sam 讲授这门课、耐心回答我关于材料的许多问题、鼓励我完成并上传这份讲义，以及对我早先草稿的评论。请注意，第 1 讲与第 16 讲很大程度上改编自 Sam 本人提供的讲义。

希望这份讲义对你有用，也希望你像我一样享受这门课。
